% kernel/engineering_identity-DE.tex

\chapter{Identität durch Übergangsänderung}
\label{chap:engineering-identity}

Der vorgeschlagene Übergang bezieht sich nun auf ein Alignment. Der Ingenieur
ändert die Länge der Eingangshalbwelle und berechnet die gekoppelten Parameter
neu. Geometrie und abgeleitete Pose ändern sich. Der Gegenstand muss sich nicht
ändern.

\section{Objektkontinuität und Zustandsnachfolge}

\begin{definition}[Engineering Object Identity]
\label{def:engineering-object-identity}
Engineering Object Identity ist die Grundlage, durch die ein Engineering
Object über seinen Ingenieurlebenszyklus unterscheidbar bleibt.
\end{definition}

Objektidentität beantwortet: \emph{Welches fortbestehende Alignment?}
Zustandsidentität beantwortet: \emph{Welcher qualifizierte Zustand oder welche
Revision dieses Alignments?} Angenommener Vorgänger und berechneter Vorschlag
können sich deshalb auf dasselbe Alignment beziehen und zugleich verschiedene
Zustände bleiben. Eine protokollierte Ableitungsrelation verbindet sie, ohne
die frühere Evidenz zu überschreiben.

Eine Parameteränderung liefert normalerweise einen möglichen Nachfolgezustand:
Alignment-Bezug, Längsordnung und nachvollziehbarer Vorgänger bleiben erhalten,
während sich die parametrisierte Konstruktion ändert. Dies ist ein erklärender
Fall, keine allgemeine Lifecycle-Regel. Die Annahme des Nachfolgezustands
verlangt weiterhin die anwendbare Autorität.

\section{Constructive Identity macht die Änderung verständlich}

Der Berlinish-Ausdruck ist nicht nur Darstellungsgeometrie. Geordnete
Komponenten, intrinsische Längsparametrisierung, Wahl der Funktionsfamilien und
Randbeziehungen halten die Konstruktion prüfbar. Der Austausch des
Klothoidenkerns gegen eine andere zulässige Familie ist daher eine konstruktive
Änderung, selbst wenn eine abgetastete Polylinie der alten sehr ähnlich bleibt.
Neuabtastung oder Transformation dieser Polylinie ist nur eine
repräsentationale Änderung, solange die Ingenieurkonstruktion nicht ebenfalls
geändert wird.

\begin{definition}[Identity Composition]
\label{def:identity-composition}
Identity Composition ist die kohärente Verbindung stabiler Identitätsaspekte
und kanonischer Relationen, die ein Objekt unterscheidbar macht, ohne Identität
auf ein repräsentationsabhängiges Attribut zu reduzieren.
\end{definition}

Bei einem Alignment trägt Constructive Identity die intrinsische Folge und
Längskonstruktion bei, durch die seine Elemente das Alignment bilden. Alignment
Identity stützt sich auf diese konstruktive Unterscheidung; sie ist weder
CRS-Kennung noch Stationsbezeichnung, Datenbankschlüssel oder erzeugte
Geometrie. Betriebliche Stationierung kann Positionen am Objekt adressieren,
ersetzt aber nicht die intrinsische Längsparametrisierung als
identitätstragende Konstruktion.

\section{Wann erzeugt eine Änderung ein anderes Objekt?}

Einige Fälle sind eindeutig. Koordinatentransformation, Serialisierung,
Rendering und Neuerzeugung abgeleiteter Geometrie ersetzen das Alignment nicht
von selbst. Die Ableitung eines editierbaren Modells unter Erhalt des
Quelldatensatzes löscht auch dessen Identität und Provenienz nicht. Am anderen
Ende kann ein ausdrücklich beauftragtes Ersatzalignment als anderes Objekt
geführt werden, selbst wenn es viel Geometrie wiederverwendet.

Dazwischen liegt eine echte Grenze. Der Austausch einer Übergangskomponente
kann als Nachfolgezustand angenommen werden oder an einem Ersatzobjekt
mitwirken, abhängig von objekttypspezifischen Gleichheitsbedingungen und dem
autorisierten Ingenieurakt. Der aktive Kernel schreibt für dieses Urteil
keinen universellen Schwellenwert vor. Geometrischer Abstand,
Softwareschlüssel oder Solver-Operation können es nicht entscheiden.

\section{Provenienz bewahrt den Weg durch die Änderung}

Der Kandidat hält Vorgänger, Berechnungsdienst und Version,
Eingangsbedingungen, gewählte Funktionsfamilien und Parameter, abgeleitete
Ergebnisse, verantwortliche Person und Zeitpunkt fest. Wird er verworfen,
bleibt diese Evidenz erhalten. Wird er angenommen, verbindet die Entscheidung
den neuen angenommenen Zustand mit seinem Vorgänger, statt Geschichte
umzuschreiben.

Identität sagt nun, welches Objekt und welchen Zustand der Vorschlag betrifft.
Sie sagt noch nicht, wo der Übergang metrisch sinnvoll ist oder ob seine
Evidenz für eine bestimmte Aufgabe genügt. Diese Abhängigkeiten begründen
Anwendbarkeit.
